% UBT-AI-PROVENANCE-BEGIN
% schema: ubt-ai-provenance/v1
% tier: C_working
% ai_assistance: disclosed
% human_review: risk-based
% editorial_responsibility: Ing. David Jaroš
% policy: ../AI_PROVENANCE.md
% notice: Working material; exhaustive human review is not claimed.
% UBT-AI-PROVENANCE-END

\documentclass[11pt,a4paper]{article}
\usepackage[T1]{fontenc}
\usepackage[utf8]{inputenc}
\usepackage[margin=2.5cm]{geometry}
\usepackage{amsmath,amssymb,amsthm}
\usepackage[hidelinks]{hyperref}
\usepackage{ubtprovenance}
\UBTTier{C}
\newtheorem{proposition}{Tvrzení}
\title{Kanonická oprava GR článku UBT:\\
hranice platnosti staršího Schwarzschildova ansatzu}
\author{Ing. David Jaro\v{s}}
\date{23. srpna 2026}

\begin{document}
\maketitle
\UBTProvenanceNotice

% BILINGUAL-UNIT: correction.scope
\section{Rozsah a závazný status}
Tato doprovodná oprava je závazná pro odstavec o Schwarzschildově statusu v
\path{papers/UBT_GR_Submission.tex}. Nemění dokázané lokální věty o tetrádě,
konexi, kontorzi ani děleném jetu. Opravuje status staršího ansatzu v
\path{canonical/geometry/biquaternionic_vacuum_solutions.tex} a skriptu
\path{tools/verify_schwarzschild_theta.py}.

Historický přímý ansatz je jako kanonické odvození neplatný. Odděleně od něj
dávají pozdější věty o efektivní akci a split-jetu závazný status obnovy
\[
 \boxed{\texttt{GAP-U2Theta: CLOSED CONDITIONALLY FOR GR RECOVERY}.}
\]
To znamená, že úplná Schwarzschildova tetráda i lapse jsou on shell v obnovené
Einsteinově infračervené větvi a mají lokální efektivní UBT reprezentaci.
Neznamená to opravu historického ansatzu ani již hotový mikroskopický výběr této
větve akcí pouze pro $\Theta$.

% BILINGUAL-UNIT: correction.psi
\section{Chybějící závislost na imaginárním čase}
Historický zdroj zapisuje
\begin{equation}
 \Theta_0(r,\tau)=e^{i\alpha(r)}X(r).
 \label{eq:legacy-ansatz}
\end{equation}
Rovnice~\eqref{eq:legacy-ansatz} neobsahuje žádnou zobrazenou závislost na
$\psi$. Proto
\begin{equation}
 \partial_\psi\Theta_0=0,
 \label{eq:psi-zero}
\end{equation}
nikoli $\partial_\psi\Theta_0=i\alpha(r)\Theta_0$. Pouhé prohlášení
$\tau=t+i\psi$ chybějící závislost nevytvoří. Historický krok
$g_{tt}=-\alpha(r)^2$ tedy z tohoto ansatzu neplyne.

% BILINGUAL-UNIT: correction.phase
\section{Vynechaná derivace radiální fáze}
Pro nekonstantní radiální fázi dává přímá derivace
\begin{equation}
 \partial_i\Theta_0=e^{i\alpha(r)}\left[
 \partial_iX+i\alpha'(r)\frac{x_i}{r}X\right].
 \label{eq:phase-derivative}
\end{equation}
Druhý člen v rovnici~\eqref{eq:phase-derivative} byl v tvrzeném zrušení
prostorové fáze vynechán. Výpočet prostorové části bez fáze proto nelze
současně použít s radiálně proměnnou fází bez nového výpočtu metriky a
smíšených členů.

% BILINGUAL-UNIT: correction.slice
\section{Překážka kanonického Lorentzova řezu}
Numerický skript kontroluje pole bez fáze s reálnými kvaterniony
\begin{equation}
 \Theta_0(r)=f(r)\mathbf1+g(r)\mathbf e_r,
 \qquad
 g(r)=r\left(1+\frac{M}{2r}\right)^2,
 \qquad
 f'(r)=\left(1+\frac{M}{2r}\right)\sqrt{\frac{2M}{r}}.
 \label{eq:legacy-spatial}
\end{equation}
Pro $M=1$, $r=2M$ a na kladné ose $x$ je skalární koeficient
$\partial_x\Theta_0$ roven
\begin{equation}
 f'(2M)=\frac54.
 \label{eq:slice-witness}
\end{equation}
Je reálný a nenulový. Kanonický Lorentzův řez naproti tomu vyžaduje skalární
koeficient ve tvaru $i e^0$ s $e^0\in\mathbb R$. Derivace používané starším
prostorovým ověřovačem jsou proto obecně mimo kanonický Lorentzův řez $W_L$.

Tentýž přesný bod také vyvrací starší pomocný předpoklad $|X|=1$, protože
\begin{equation}
 g(2M)=\frac{25}{8},
 \qquad |X|^2=f^2+g^2\geq\frac{625}{64}>1.
 \label{eq:norm-witness}
\end{equation}

% BILINGUAL-UNIT: correction.verdict
\section{Opravený závěr}
Identita
\begin{equation}
 f'(r)^2+g'(r)^2=\left(1+\frac{M}{2r}\right)^4
 \label{eq:spatial-identity}
\end{equation}
platí a starší skript reprodukuje kladnou prostorovou kvadratickou formu
$g_{ij}=\Psi^4\delta_{ij}$. Tato omezená identita není kanonickým
Schwarzschildovým řešením z jediného $\Theta$, neurčuje lapse a není on-shell
výsledkem.

Skutečné uzavření obnovy používá jiné, pozdější výsledky. Podmíněná
infračervená efektivní akce dává Einsteinovu--$\Lambda$ rovnici. Pro
$\Lambda=0$ je úplná Schwarzschildova tetráda včetně lapse vakuovým
Einsteinovým řešením. Věta o split-jet pravé inverzi reprezentuje každou
regulární lokální tetrádu při zachování fyzické Levi--Civitovy konexe a
split-jet multiplikátor je na slupce nulový. Schwarzschildovo řešení tedy patří
do podmíněné efektivní UBT GR větve bez použití neplatného historického ansatzu.
Tato kompozice je explicitně dokázána v
\path{canonical/gr_closure/gr_recovery_completion.cs.tex} a zaznamenána v
\path{canonical/gr_closure/gr_recovery_status.yaml}.

Silnější mikroskopický problém---odvodit úplný efektivní metrický selektor
přímo z finalizované akce pouze pro $\Theta$ a globálně pokračovat reprezentaci
přes všechny nulové/horizontové patche---zůstává otevřeným problémem
fundamentálního dokončení, nikoli blokátorem podmíněné lokální obnovy GR.

% BILINGUAL-UNIT: correction.verification
\section{Záznam ověření}
Analytickou opravu nezávisle kontrolují
\path{tools/verify_schwarzschild_theta.py} a nativní SageMath v
\path{tools/verify_gr_core_sage.py}. Uzavření obnovy skládá dříve dokázané věty
o split-jetu a efektivní Einsteinově větvi. Úplný strojově čitelný záznam je v
\path{papers/UBT_GR_Submission.verification.yaml}.

\textbf{Status Lean: LEAN-PENDING.} Lean a \texttt{lake} nebyly v prostředí
auditu dostupné a žádný existující zdroj Lean tyto GR výroky neformalizuje.
Kontroly CAS ověřují pouze výše zapsané rovnice. Nedokazují mikroskopický původ
akce, globální pokračování ani fyzickou pravdivost UBT.

\end{document}
